\documentclass[11pt]{article}

% ===================== PACKAGES =====================
\usepackage[a4paper,margin=1in]{geometry}
\usepackage{amsmath,amssymb}
\usepackage{enumitem}
\usepackage{fancyhdr}
\usepackage{xcolor}

% ===================== HEADER & FOOTER =====================
\pagestyle{fancy}
\fancyhf{}
\lhead{CSE 400: Fundamentals of Probability in Computing}
\rhead{Lecture 11}
\cfoot{\thepage}

% ===================== TITLE =====================
\title{
    \normalsize School of Engineering and Applied Science (SEAS), Ahmedabad University \\
    \vspace{0.2cm}
    \textbf{CSE 400: Fundamentals of Probability in Computing}\\
    \Large Lecture Scribe Submission
}
\author{}
\date{}

\begin{document}
\maketitle

\vspace{-2cm}
\begin{center}
    \begin{tabular}{ll}
        \textbf{Name:} & Ishika Patel \\[1.5ex]
        \textbf{Enrollment No:} & AU2440088 \\[1.5ex]
        \textbf{Group:} & s1\_g1\_fin \\[1.5ex]
        \textbf{Lecture:} & 11 \\[1.5ex]
        \textbf{Instructor:} & Dhaval Patel, PhD \\[1.5ex]
        \textbf{Date:} & February 10, 2026
    \end{tabular}
\end{center}

\hrule
\vspace{0.5cm}

\section{Transformation of Random Variables}

\subsection{Objective}

Learning of transformation techniques for random variables.

\subsection{Single Random Variable Transformation}

Let $X$ be a random variable whose PDF $f_X(x)$ is known.

Let
\[
Y = g(X)
\]

We aim to find the CDF and PDF of $Y$.

\subsubsection*{Step–V: CDF of $Y$}

Since
\[
F_Y(y) = \Pr(Y \le y)
\]

\[
F_Y(y) = \Pr(g(X) \le y)
\]

This probability is expressed in terms of $X$ and evaluated using the known PDF of $X$.

\subsubsection*{PDF of $Y$}

For a monotonic transformation,
\[
f_Y(y) = f_X(x) \left| \frac{dx}{dy} \right|
\]

evaluated at the value of $x$ such that
\[
y = g(x)
\]

If the mapping is not one-to-one, contributions from all valid roots are added.

If $y$ is outside the transformed range, then
\[
f_Y(y) = 0
\]

\section{Function of Two Random Variables}

Joint transformations and derived distributions.

Let
\[
Z = X + Y
\]

Then the CDF of $Z$ is
\[
F_Z(z) = \Pr(Z \le z)
\]

\[
= \Pr(X + Y \le z)
\]

This probability is evaluated over the joint distribution of $X$ and $Y$.

\section{Illustrative Examples}

\subsection{Example 1}

Let
\[
X \sim \text{Uniform}(-1,1)
\]

\[
f_X(x) =
\begin{cases}
\frac{1}{2}, & -1 < x < 1 \\
0, & \text{otherwise}
\end{cases}
\]

Let
\[
Y = \sin\left(\frac{\pi}{2} X\right)
\]

Find the PDF of $Y$.

\subsubsection*{Step 1: Express the transformation}

\[
Y = \sin\left(\frac{\pi}{2} X\right)
\]

\subsubsection*{Step 2: Find the inverse transformation}

\[
X = \frac{2}{\pi} \sin^{-1}(Y)
\]

\subsubsection*{Step 3: Compute derivative}

\[
\frac{dX}{dY} = \frac{2}{\pi} \cdot \frac{1}{\sqrt{1 - Y^2}}
\]

\subsubsection*{Step 4: Apply transformation formula}

\[
f_Y(y) = f_X(x) \left| \frac{dX}{dY} \right|
\]

Substituting,

\[
f_Y(y) =
\frac{1}{2}
\cdot
\frac{2}{\pi}
\cdot
\frac{1}{\sqrt{1-y^2}}
\]

\[
f_Y(y) =
\frac{1}{\pi \sqrt{1-y^2}},
\quad -1 < y < 1
\]

\[
f_Y(y) = 0 \quad \text{otherwise}
\]

\subsection{Example 2}

Let
\[
Z = X + Y
\]

Then
\[
F_Z(z) = \Pr(Z \le z)
\]

\[
= \Pr(X + Y \le z)
\]

This is evaluated using the joint distribution of $X$ and $Y$.

\section{Summary of Lecture}

Transformation of a random variable is performed using the CDF method or PDF transformation formula.

For monotonic $g(x)$:
\[
f_Y(y) = f_X(x) \left| \frac{dx}{dy} \right|
\]

For two random variables:
\[
F_Z(z) = \Pr(X + Y \le z)
\]

Example derived:
\[
f_Y(y) =
\frac{1}{\pi \sqrt{1-y^2}},
\quad -1 < y < 1
\]

\bigskip
\hrule
\vspace{0.2cm}
\begin{center}
    \small \textit{End of Lecture 11}
\end{center}

\end{document}
